La complétion cubique et le transport apériodique précisent deux sens
de la conservation qui s'articulent dans le développement. La première conserve
l'unité par projection de mesures normalisées et distingue le cube
initial, la complétion perforée de \(999\) éléments et la restitution
du sommet. Le second conserve le registre de trajectoire durant le
retour de phase. La monodromie et le modèle mathématique de cristal temporel
apériodique expriment ce second mécanisme au moyen d'applications et
d'observables, et non seulement par une image hélicoïdale. L'exponentielle
tritique, quant à elle, unifie les réalisations elliptique, parabolique et
hyperbolique sans supprimer leurs relations quadratiques respectives.

La constante de Planck réduite, \(\hbar\), occupe une position explicite dans cette architecture. Le quotient déterminantal d'ordre \(16\) agit dans la norme qui détermine la décade \(-34\) ; le caractère \(H_5\) et sa continuation régionale déterminent les deux sections d'action dont le quotient entre dans la coordonnée électronique. La comparaison de la section précédente avec \(h/(2\pi)\) montre une différence relative de l'ordre de \(10^{-15}\), et la section retournée conserve la compensation qu'utilise l'électron. L'échelle d'action relie ainsi la norme déterminantale et le caractère régional au quotient électronique : la compensation appartient à la construction exacte, tandis que la proximité de la section précédente avec \(h/(2\pi)\) est une comparaison numérique et non une égalité exacte.
